\bookchapter{Глава 13: Рекомендации приложений-инструментов}

\chapterintro{В этой главе собраны инструменты для продуктивности, здоровья и поездок, а также полезные системные и межустройственные возможности: очистка и защита от спама (Qingzhuo и др.), встроенные переключатели блокировки в самом телефоне, предупреждения о погоде и землетрясениях; удалённое управление компьютером (Sunlogin/ToDesk), ИИ-офис вроде Trae, YoMate, быстрая передача файлов по локальной сети (LocalSend и др.), общий буфер обмена между устройствами, одноразовые коды Aegis/Microsoft Authenticator, планировщик TickTick и сканирование бумажных документов.}

\booksection{Инструменты продуктивности: время и внимание}

Инструменты продуктивности — «скелет» цифровой жизни: помогают привести в порядок разрозненные мысли, задачи и материалы. Ниже — рекомендации по сценариям.

\booksubsection{Списки задач (управление задачами)}

\begin{itemize}
  \item TickTick. Для кого: для тех, кому нужно вести в одном месте задачи и расписание работы, личной жизни и учёбы; одна из ключевых рекомендаций этого раздела. Главные плюсы: синхронизация между телефоном/компьютером/планшетом/Web; списки + календарь + помодоро-таймер — три в одном; ввод на естественном языке («встреча завтра в 15:00», «отчёт каждый понедельник») с автоматическим распознаванием времени и правил повтора; поддержка подзадач, тегов, приоритетов, вложений и совместных списков. Подходит для комплексного управления по схеме «расписание + задачи + ежедневный обзор».
  \item Как удобнее пользоваться TickTick: Инбокс: сначала всё скидывайте в «Сегодня/Инбокс», вечером раскладывайте по спискам «Работа», «Дом», «Учёба».
  \item Календарный вид: перетаскивайте встречи и дедлайны в календарь, чтобы видеть конфликты нагляднее, чем в простом списке.
  \item Повторяющиеся задачи: лекарства, платежи, тренировки задавайте правилами, чтобы не вводить каждый раз вручную.
  \item Помидорка: запускайте фокус прямо из задачи и смотрите, сколько на этой неделе вы реально сделали.
  \item Связь с функциями телефона: голосовой помощник «Напомни…» отправит в системные напоминания; на важные дедлайны можно дублировать в системный календарь.
\end{itemize}

Todoist Для кого: любителям минимализма и кроссплатформенности. Ключевые плюсы: спокойный интерфейс, чёткая иерархия (задача подзадача подподзадача), карма-баллы для привычек.

Microsoft To Do Для кого: офисным работникам и пользователям Outlook. Ключевые плюсы: полностью бесплатно, без рекламы, функция «Мой день» очищается каждый день и снижает нагрузку, глубокая интеграция с Microsoft.

Системные заметки/напоминания Для кого: лёгким пользователям, не желающим ставить сторонние приложения. Ключевые плюсы: нативный опыт, нулевая задержка, добавление через Siri/Алису, работа с виджетами.

\booksubsection{Помидорка (тренировка концентрации)}

\begin{itemize}
  \item Forest Для кого: тем, кто не может удержаться от телефона. Ключевые плюсы: во время фокуса растёт дерево; если выйти из приложения раньше — оно засохнет; положительная обратная связь через «виртуальную жизнь», можно сажать деревья с друзьями.
  \item Tide. Для кого: для тех, кому нужен белый шум для концентрации или засыпания. Главные плюсы: три в одном — природные звуки (дождь, море, лес), помодоро-таймер и медитация; минималистичный, умиротворяющий интерфейс.
  \item Принцип: классика — 25 минут фокуса + 5 минут отдыха; можно подстроить под себя 50/10 или 90/20.
\end{itemize}

\booksubsection{Календарь (визуализация времени)}

\begin{itemize}
  \item Системный календарь Для кого: всем. Ключевые плюсы: самые стабильные push-уведомления, низкое энергопотребление, добавление через Siri/Алису.
  \item Google Calendar Для кого: тем, у кого много устройств и нужно делиться расписанием. Ключевые плюсы: функция «Цели» сама подбирает время для тренировок/учёбы, умное распознавание событий из Gmail.
\end{itemize}

\booksubsection{Заметки (систематизация знаний)}

\begin{itemize}
  \item Notion Для кого: тем, кто хочет вести заметки, задачи, базы данных и вики в одном месте; студентам, фрилансерам, командам. Ключевые плюсы: универсальная база знаний, блочное редактирование, можно собрать личную страницу, список чтения, доску проектов, богатая библиотека шаблонов.
  \item Obsidian Для кого: глубоким управленцам знаниями, исследователям, авторам. Ключевые плюсы: локальное хранение Markdown, двусторонние ссылки и граф связей — заметки превращаются во «второй мозг», данные полностью ваши.
  \item flomo Для кого: кто хочет фиксировать мысли, но не писать длинные тексты. Ключевые плюсы: карточные заметки, запись идей как в соцсети, метки \# и возврат по таймлайну.
  \item Заметки Apple Для кого: пользователям iPhone/Mac с лёгкими потребностями. Ключевые плюсы: быстрый запуск, сканирование документов, рисование, списки, синхронизация iCloud, бесплатно и глубоко интегрировано с системой.
  \item Evernote. Для кого: для опытных пользователей, которым нужно много вырезать из веба и архивировать материалы. Главные плюсы: мощное расширение для веб-клиппинга, распознавание текста на картинках (OCR), проверенный временем стабильный сервис, но за расширенную версию придётся платить.
\end{itemize}

\booksubsection{PDF-инструменты}

\begin{itemize}
  \item PDF Expert Для кого: пользователям iPad/iPhone, бизнесменам с договорами и студентам. Ключевые плюсы: признан лучшим PDF-редактором на iOS, удобная разметка, подписи, объединение/разделение, заполнение форм.
  \item WPS Для кого: тем, кто работает с Word/Excel/PPT/PDF. Ключевые плюсы: единый офисный пакет, хорошая совместимость с китайскими облаками и файлами WeChat, бесплатной версии хватает на каждый день.
  \item Adobe Acrobat Для кого: кому важна строгая спецификация PDF (юриспруденция, дизайн, издательское дело). Ключевые плюсы: официальный стандарт, безошибочное преобразование в PDF/A, профессиональная синхронизация и электронная подпись.
\end{itemize}

\booksubsection{Сканирование (оцифровка бумаг)}

Договоры, чеки, домашние задания, записи с доски, копии паспорта… Камера телефона уже заменяет большинство сканеров. Основной процесс: снять чётко автообрезка и выравнивание экспорт в PDF/картинку при желании OCR.

\begin{itemize}
  \item CamScanner. Для кого: для деловых людей и студентов, которым часто нужно сканировать договоры, чеки и документы. Главные плюсы: автоматическое распознавание границ, исправление перспективы, OCR-распознавание текста; умеет создавать многостраничные PDF и синхронизировать их в облаке.
  \item Microsoft Lens Для кого: студентам и тем, кто хочет сканировать доску/страницы в Word. Ключевые плюсы: бесплатно, без рекламы, экспорт в OneNote/Word в один клик, высокая точность распознавания.
  \item Сканирование в Заметках Apple Для кого: пользователям iPhone с лёгкими задачами. Ключевые плюсы: не нужно ставить приложение — в Заметках нажмите иконку камеры, автоудаление теней, подписи, бесплатно и удобно.
  \item Baimiao / веб-версия CamScanner / режим «Скан документа» в камере бренда: у Huawei, Xiaomi, OPPO и vivo режим «Документ/Скан» по большей части встроен в системную камеру; на Android подойдут также Google Camera или системная функция сканирования. Годится, чтобы наспех отсканировать страницу договора.
  \item Советы по съёмке: ровный свет, телефон параллельно бумаге, наводите на углы; многостраничные документы снимайте по порядку и объединяйте в PDF; для чувствительных документов добавляйте водяные знаки и следите за местом хранения, после использования очищайте «Недавно удалённые».
\end{itemize}

\booksection{Здоровье и спорт}

Приложения для здоровья — самый недооценённый «личный помощник» телефона: они постоянно следят за сном, пульсом и активностью и лучше отражают долгосрочное состояние, чем редкие визиты к врачу.

\booksubsection{Фитнес}

\begin{itemize}
  \item Keep Для кого: тренирующимся дома, новичкам, тем, кто хочет похудеть и подтянуться. Ключевые плюсы: множество бесплатных курсов от HIIT до йоги, ИИ составляет план, живое сообщество.
  \item Nike Training Club (NTC) Для кого: тем, кто хочет профессионального тренера и системных тренировок. Ключевые плюсы: официальный продукт Nike, все курсы бесплатные, движения показывают профессиональные тренеры, удобная классификация по мышцам и инвентарю.
  \item Apple Fitness Для кого: владельцам Apple Watch и тем, кого мотивируют «кольца активности». Ключевые плюсы: три ежедневных кольца (движение/тренировка/стоя), наглядные цели, бесшовная связь с пульсом часов.
\end{itemize}

\booksubsection{Отслеживание сна}

\begin{itemize}
  \item AutoSleep Для кого: владельцам Apple Watch, кто хочет знать долю глубокого и лёгкого сна. Ключевые плюсы: автоматическое отслеживание с часами на руке, без ручного включения, оценка качества сна и советы по режиму.
  \item Sleep Cycle Для кого: тем, кто спит без часов и хочет следить телефоном. Ключевые плюсы: анализ храпа и движений через микрофон, умное пробуждение в лёгкой фазе — вставать проще.
  \item Xiaomi Fitness Для кого: владельцам браслетов и часов Xiaomi. Ключевые плюсы: бесплатно и глубоко связано с браслетом Xiaomi, пульс/кислород/сон в одном трекере, отличное соотношение цены и качества.
\end{itemize}

\booksubsection{Медитация и расслабление}

\begin{itemize}
  \item Calm Для кого: занятым офисным работникам и тем, кому трудно засыпать. Ключевые плюсы: ежедневные медитации «Daily Calm» + сказки на ночь (голосами знаменитостей), красивый интерфейс, много англоязычного контента.
  \item Tide. Для кого: для китайскоязычных пользователей, которым нужен белый шум для концентрации или засыпания. Главные плюсы: китайский продукт премиум-качества — природные звуки + курсы медитации + помодоро-таймер; спокойные подсказки на родном языке звучат теплее.
  \item SnailSleep. Для кого: для тех, кому нужны звуки для засыпания и запись разговоров во сне/храпа. Главные плюсы: огромная библиотека белого шума со слоями (дождь + вентилятор + костёр), таймер отключения перед сном; запись сна (разговоров во сне) — и забавно, и полезно.
\end{itemize}

\booksubsection{Женское здоровье}

\begin{itemize}
  \item Flo Для кого: кто ведёт календарь менструаций, планирует беременность или следит за женским здоровьем. Ключевые плюсы: ИИ точно предсказывает следующую менструацию и овуляцию, анонимные вопросы сообщества, изящный интерфейс.
  \item Meiyou. Для кого: для женщин в Китае и тех, кто хочет китайскоязычное сообщество и научно-популярные материалы. Главные плюсы: богатый локализованный контент — знания подбираются по этапам: цикл / планирование беременности / беременность / уход за малышом.
\end{itemize}

\booksubsection{Напоминания о лекарствах}

\begin{itemize}
  \item «Помощник по приёму лекарств». Для кого: для пациентов с хроническими болезнями, принимающих лекарства постоянно, и для взрослых детей, ухаживающих за пожилыми родителями. Главные плюсы: распознавание лекарства по штрихкоду, автоматическое составление списка приёма, вовремя всплывающие напоминания; записи о приёме можно отправить родственникам.
\end{itemize}

\booksubsection{Центр данных}

\begin{itemize}
  \item Apple Health / Google Fit Для кого: всем пользователям смартфонов. Ключевые плюсы: собирают данные с телефона, браслета и сторонних приложений (шаги, пульс, сон, менструации и т. д.) в единую «карту здоровья».
\end{itemize}

\booksection{Путешествия и поездки}

В поездке несколько удобных приложений избавят от множества проблем. Ниже — рекомендации по всему пути.

\booksubsection{Авиа и поезда}

\begin{itemize}
  \item Umetrip. Для кого: для деловых людей, которые часто летают. Главные плюсы: официальные данные гражданской авиации — самые надёжные сведения о статусе рейса, смене выхода на посадку и предупреждениях о задержках; можно посмотреть модель самолёта и схему салона.
  \item 12306 Для кого: всем, кто ездит на поездах по Китаю. Ключевые плюсы: официальный и единственный надёжный канал покупки билетов, лист ожидания (автопокупка при возврате), удобные студенческие и групповые билеты.
  \item VariFlight. Для кого: для тех, кто чувствителен к пунктуальности рейсов и встречает/провожает в аэропорту. Главные плюсы: точный прогноз задержек; можно поделиться поездкой с семьёй — близкие в реальном времени видят, где вы и когда прибудете.
\end{itemize}

\booksubsection{Карты и навигация}

\begin{itemize}
  \item Amap (Gaode Maps). Для кого: для пользователей в китайских городах, водителей и тех, кто ездит на такси. Главные плюсы: точная дорожная обстановка в реальном времени и умный объезд пробок; такси и совместные поездки — всё в одном месте; богатый выбор вариантов пересадок на общественном транспорте.
  \item Baidu Maps Для кого: кому нужна информация о местах и навигация внутри помещений. Ключевые плюсы: панорамы/AR-навигация, детальные карты ТЦ и достопримечательностей, много голосов.
  \item Google Maps Для кого: выезжающим за границу. Ключевые плюсы: лучшее глобальное покрытие, честные местные отзывы, Street View и офлайн-карты спасают в незнакомой стране.
\end{itemize}

\booksubsection{Перевод}

\begin{itemize}
  \item Google Translate Для кого: путешественникам и изучающим языки. Ключевые плюсы: режим «Разговор» с двухсторонним переводом в реальном времени, камера для мгновенного перевода текста — незаменимо в поездке, можно скачать офлайн-пакеты.
\end{itemize}

\booksubsection{Конвертация валют}

\begin{itemize}
  \item «Минималистичный курс валют». Для кого: для тех, кто ездит за границу, заказывает с зарубежных сайтов и делает трансграничные покупки. Главные плюсы: предельно чистый интерфейс, поддержка более 160 валют, курсы обновляются в реальном времени, несколько валют можно сравнивать одновременно, а долгое нажатие на число мгновенно пересчитывает сумму.
\end{itemize}

\booksubsection{Погода}

\begin{itemize}
  \item Системное приложение «Погода» / China Weather. Для кого: для всех; первый выбор для прогноза и официальных метеопредупреждений. Главные плюсы: погодные приложения производителей и официальные каналы вроде «China Weather» связаны с предупреждениями метеослужб — уведомления о ливнях, тайфунах, жаре приходят баннером или пушем, и доверять им можно.
  \item Caiyun Weather. Для кого: для тех, кому нужен прогноз дождя с точностью до минуты. Главные плюсы: точность ответа на вопрос «будет ли дождь в ближайшие два часа» — одна из лучших в отрасли; наглядная динамическая карта дождевых фронтов; приложение само предупредит, прежде чем начнётся дождь.
  \item Carrot Weather Для кого: любителям погоды, ценящим юмор и виджеты. Ключевые плюсы: ироничный голос ИИ, множество настраиваемых виджетов, профессиональная радарная карта; платная версия мощнее.
  \item Важно о предупреждениях: приложению погоды нужно разрешить уведомления, чтобы вы получали оповещения; уровни — синий, жёлтый, оранжевый, красный; красный — максимальный риск, меняйте планы поездок. Подробнее о включении — в разделе 7 этой главы «Предупреждения о погоде / стихийных бедствиях / землетрясениях».
\end{itemize}

\booksubsection{Гиды для путешествий}

\begin{itemize}
  \item Mafengwo. Для кого: для тех, кто любит подробные путеводители, автопутешествия и глубокие поездки. Главные плюсы: богатый авторский контент — путевые заметки пользователей, детальные описания развлечений, маршрутов и «граблей» по каждому направлению; гайды можно скачивать для офлайна.
  \item Xiaohongshu (RED). Для кого: для молодых путешественников, которые хотят реальные впечатления и ищут малоизвестные места и вкусную еду. Главные плюсы: наглядный формат из фото, текстов и коротких видео; подборки вроде «Город XX за 3 дня» обновляются быстро и всегда актуальны; удобный поиск.
\end{itemize}

\booksection{Введение в автоматизацию}

Автоматизация — рычаг «настроил один раз, экономишь время долго»: десятки повторяющихся действий сжимаются до нуля. Здесь только концепция; подробные сценарии и уроки — в главе 23: «Программное мышление и введение в автоматизацию».

\begin{itemize}
  \item «Быстрые команды» iOS: официальный инструмент автоматизации Apple — без программирования можно одним касанием выполнить связку «дома включить Wi-Fi + уменьшить яркость + прочитать погоду».
  \item Автоматизация Android: системные сценарии (например, «Автозадачи» Xiaomi, «Контекстный помощник» Huawei) освоить проще всего; Tasker мощнее, но с крутой кривой обучения — для гиков.
  \item Главная идея: сначала определите «действия, которые повторяются каждый день», затем решите, можно ли заменить их «одним триггером + несколькими действиями». Подробные примеры и продвинутые приёмы — в главе 23.
\end{itemize}

\booksection{Браузеры и поисковики: ваш вход в интернет}

Браузер — окно в интернет, поисковик — компас в океане информации. Умное использование сэкономит немало времени.

\booksubsection{Как работает браузер:}

Когда вы вводите адрес, браузер: (1) запрашивает у DNS IP-адрес домена; (2) отправляет HTTP/HTTPS-запрос на этот IP; (3) сервер возвращает код страницы (HTML, CSS, JavaScript); (4) браузер разбирает код и рисует привычную вам страницу.

\booksubsection{Рекомендуемые браузеры:}

\begin{itemize}
  \item Microsoft Edge: на ядре Chromium, быстрый и совместимый. Встроенный помощник Copilot умеет резюмировать страницы, переводить и вести диалог. Синхронизация между устройствами — лично я настоятельно рекомендую его как основной браузер.
  \item Safari: браузер по умолчанию на устройствах Apple, глубоко интегрирован с системой, экономит батарею. Хорошая приватность, но меньше расширений.
  \item Chrome: самое высокое использование в мире, богатые расширения. Но много ест память и на телефоне сильнее расходует батарею.
\end{itemize}

\booksubsection{Выбор поисковика:}

\begin{itemize}
  \item Bing: поисковик Microsoft, тесно интегрирован с Edge. ИИ-резюме сразу даёт ответ, а не только ссылки — поиск эффективнее. Хороший поиск по картинкам.
  \item Yandex: крупнейший поисковик России, сильный поиск по картинкам и обратный поиск. Отлично подходит для поиска похожих изображений и распознавания объектов.
  \item Baidu: основной поисковик в Китае, но много рекламы — нужна критичность.
  \item Google: самый точный поисковик в мире, но в Китае требует особого доступа к сети.
\end{itemize}

\booksubsection{Как искать информацию эффективно:}

\begin{itemize}
  \item Используйте ИИ-поиск: резюме Bing и Copilot в Edge сразу отвечают на вопросы и выделяют главное из массы результатов — куда быстрее, чем открывать ссылки по одной.
  \item Распознавайте рекламу и фейки: записи с пометкой «Реклама» или «Спонсор» — платное продвижение, а не органическая выдача. Проверяйте источник: предпочитайте официальные сайты и известные СМИ.
  \item Правильно скачивайте приложения: введите «название приложения + официальный сайт», например «WeChat официальный сайт», и открывайте результаты с официальным доменом. Не ставьте ПО со сторонних загрузчиков — они часто тянут мусор.
  \item Как отличить официальный сайт: Windows 11 нужно качать с microsoft.com, Chrome — с google.com. Совпадение домена с брендом — самый простой способ проверки.
\end{itemize}

\booksubsection{База знаний о доменах:}

Домен — «адрес» сайта. Понимание суффиксов помогает оценить природу и происхождение сайта:

\begin{itemize}
  \item Частые суффиксы: .com — коммерческая организация, самый распространённый
  \item .cn — национальный домен Китая
  \item .org — некоммерческая организация
  \item .net — сетевые сервисы
  \item .edu — образовательные учреждения
  \item .gov — государственные органы
  \item .xyz/.tech/.info — новые универсальные домены, низкий порог входа, будьте осторожны
\end{itemize}

Страновые/региональные домены: .jp — Япония

.kr — Корея

.uk — Великобритания

.de — Германия

.tw — Тайвань

.hk — Гонконг

О безопасности: сайты с .gov и .edu обычно надёжнее — порог получения домена выше. Дешёвые домены вроде .xyz и .top часто используют фишинговые сайты, будьте внимательны.

\booksubsection{Блокировка рекламы:}

\begin{itemize}
  \item AdGuard Для кого: всем, кого мешают всплывающие окна, баннеры и реклама в приложениях. Ключевые плюсы: системная блокировка рекламы (не только в браузере), фильтрация баннеров в приложениях и трекинга на заставках, экономия трафика и ускорение загрузки; платная версия — для всех устройств.
\end{itemize}

\booksubsection{Приватные браузеры:}

\begin{itemize}
  \item Brave Для кого: тем, кто ценит приватность и не хочет рекламного трекинга. Ключевые плюсы: на ядре Chromium, из коробки блокирует рекламу и трекеры, есть режим Tor, быстрее Chrome и меньше ест память.
  \item DuckDuckGo Для кого: тем, кто не хочет, чтобы после одного запроса за ним гналась реклама по всему интернету. Ключевые плюсы: собственный поисковик (не хранит историю), мгновенная очистка всех следов («кнопка сжигания»), единый опыт на разных платформах.
\end{itemize}

\booksubsection{Сравнение браузеров:}

\begin{center}
\begin{tabularx}{\linewidth}{|X|X|X|}
\hline
\textbf{Браузер} & \textbf{Ключевые плюсы} & \textbf{Для кого} \tabularnewline
\hline
Microsoft Edge & Ядро Chromium + встроенный Copilot AI: резюме страниц, перевод, диалог & Офисным работникам, ценящим эффективность, и активным пользователям ИИ (рекомендуется как основной) \tabularnewline
Safari & Глубокая интеграция с экосистемой Apple, экономия батареи, сильная приватность & Пользователям iPhone/Mac, важным автономность и связь с системой \tabularnewline
Chrome & Самое высокое использование в мире, самая богатая экосистема расширений & Продвинутым пользователям с множеством расширений и тем, кому нужна синхронизация устройств \tabularnewline
Brave & Из коробки блокирует рекламу и трекеры, работает быстрее & Тем, кто ценит приватность и хочет блокировать рекламу без плагинов \tabularnewline
DuckDuckGo & Собственный поисковик не хранит историю, следы чистятся в один клик & Чувствительным к приватности пользователям, не желающим персонализированного трекинга \tabularnewline
\hline
\end{tabularx}
\end{center}

\begin{tcbpractice}{Попробуйте создать автоматизацию «Расслабление перед сном»}
\begin{enumerate}
  \item Создайте триггер: определённое время (например, 22:30).
  \item Добавьте действия: уменьшить громкость > включить «Не беспокоить» > воспроизвести расслабляющий плейлист > уменьшить яркость.
  \item Сохраните и протестируйте.
\end{enumerate}

\end{tcbpractice}

\booksection{Практические системные инструменты: чистка, защита от спама и безопасность}

Помимо приложений продуктивности и здоровья, на повседневный опыт сильнее всего влияют три вещи: «меньше мешать», «хватает места» и «мошенники не проходят». Сначала разберём полезные утилиты, затем подробно — встроенную блокировку спама: многие пользуются телефонами годами и ни разу не открывали эти переключатели.

\booksubsection{Очистка системы и оптимизация хранения}

\begin{itemize}
  \item Qingzhuo. Для кого: для пользователей Android, которые хотят легко чистить кэш, остаточные файлы и пустые папки, но не желают мириться с навязчивым «семейством приложений Phone Manager». Главные плюсы: маленький размер, без рекламных неудобств; глубокое сканирование остатков приложений, поиск крупных файлов, анализ занятого места; сдержаннее коммерческих «ускорителей в один клик» и отлично подходит для регулярного обслуживания тем, кто не боится разбираться сам. Внимание: перед очисткой проверяйте пути; в WeChat, галерее и папках документов не выбирайте всё подряд вслепую; системный кэш по возможности чистите встроенными средствами бренда, а сторонние утилиты используйте лишь как дополнение.
  \item Встроенная очистка хранения Для кого: всем, особенно пожилым и тем, кто не любит сложности. Ключевые плюсы: Android [Настройки] > [Память] > [Очистка и ускорение]; iOS [Настройки] > [Основные] > [Хранилище iPhone]. Система лучше знает структуру кэша своих приложений, риск случайного удаления минимален — это первый выбор.
  \item Управление хранилищем WeChat/QQ Для кого: активным переписчикам. Ключевые плюсы: WeChat [Я] > [Настройки] > [Основные] > [Место хранения] — точечная очистка по чатам и типам файлов; безопаснее внешних «чистильщиков WeChat».
  \item Принцип: если есть штатное средство системы/приложения — не ставьте стороннее; если можно выбрать по пути/типу — не нажимайте «очистить всё». Цель очистки — освободить место, а не стереть важные данные.
\end{itemize}

\booksubsection{Защита от спама и безопасность}

\begin{itemize}
  \item National Anti-Fraud Center Для кого: всем, особенно пожилым, работникам финансов и тем, кто часто покупает онлайн и берёт незнакомые звонки. Ключевые плюсы: официальный продукт; предупреждения о звонках, риск SMS, проверка подозрительных ссылок; связь с базой полиции — самый официальный уровень «трёх линий защиты».
  \item Центр безопасности бренда / Менеджер телефона Для кого: пользователям соответствующих брендов. Ключевые плюсы: Huawei Phone Manager, Xiaomi Security Center, OPPO Phone Manager, vivo Phone Manager и др. — блокировка спама, антивирус, управление правами, защита платежей; обычно бесплатно и лучше ладят с системой. Второй коммерческий антивирус ставить не стоит — будут конфликты прав.
  \item Антиспам-сервисы операторов. Для кого: для тех, кого бомбардируют частыми звонками и ИИ-голосовыми вызовами. Главные плюсы: China Mobile — отправьте «KTFSR» на 10086; China Telecom — подпишитесь на «Tianyi Anti-Harassment»; China Unicom — настройте в приложении «Unicom Phone Manager». Фильтрация происходит на стороне оператора, и часть спам-номеров отсекается, даже когда телефон молчит.
  \item AdGuard Для кого: всем, кого мешают всплывающие окна, баннеры и реклама в приложениях. Ключевые плюсы: системная блокировка рекламы (не только в браузере), фильтрация баннеров и трекинга на заставках, экономия трафика и ускорение загрузки; платная версия — для всех устройств.
  \item Осторожно с «пакетами очистки и блокировки»: некоторые сторонние менеджеры постоянно работают в фоне и снова и снова запрашивают доступ к контактам и SMS; блокировка не лучше системной, а расход батареи и риски приватности растут. Сначала включайте штатные функции, затем официальное анти-фрод приложение.
\end{itemize}

\textbf{Принципы выбора утилит}\par

\begin{itemize}
  \item Сначала система: блокировка, очистка, права — включайте встроенные переключатели.
  \item Официальные дополнения: National Anti-Fraud Center и защита оператора — официальные каналы, можно совмещать.
  \item Сторонние — с ограничением: одной программы на функцию достаточно, чтобы не конфликтовали права и фон.
  \item Чем меньше, тем лучше: больше утилит — больше уведомлений и прав, выше шанс, что мешать будет сам инструмент.
\end{itemize}

\booksubsection{Встроенная блокировка спама: как включить по брендам}

Ниже — входы по брендам. Названия могут немного отличаться в разных версиях системы — ищите в настройках «спам», «блокировка», «неизвестные вызовы». Дополняет главу 8 «Борьба со спам-звонками»: здесь — где включать, там — маркировка номеров и апелляции.

\begin{center}
\begin{tabularx}{\linewidth}{|X|X|X|}
\hline
\textbf{Платформа} & \textbf{Путь входа} & \textbf{Рекомендуемые опции} \tabularnewline
\hline
iPhone / iOS & {}[Настройки] > [Телефон] / [Сообщения] & Неизвестные вызовы беззвучно; заблокированные контакты; «Фильтровать неизвестных отправителей» в Сообщениях \tabularnewline
Xiaomi / Redmi & Телефон > Настройки > Блокировка спама; или Центр безопасности & Умная блокировка, отсечение звонков «один гудок», фильтрация SMS по ключевым словам \tabularnewline
Huawei / Honor & Телефон > Блокировка спама; Менеджер телефона > Блокировка спама & Идентификация звонящих, чёрный список, умная блокировка SMS \tabularnewline
OPPO / realme / OnePlus & Телефон > Блокировка спама; или Устройство > Блокировка спама & Блокировка звонков, распознавание мошенничества, правила фильтрации SMS \tabularnewline
vivo / iQOO & Телефон > Блокировка спам-звонков; i-менеджер & Блокировка спама, предупреждения о мошенничестве, фильтрация сообщений \tabularnewline
Samsung & Телефон > Настройки > Блокировка номеров / Определение номера и спам & Блокировка неизвестных номеров, распознавание спам-звонков (поддержка зависит от региона) \tabularnewline
\hline
\end{tabularx}
\end{center}

\begin{tcbinfo}{Совмещение трёх уровней защиты}
Рекомендуется связка «Системная блокировка + защита оператора + National Anti-Fraud Center»: система занимается маркировкой и чёрным списком, оператор отсекает частые обзвоны на своей стороне, анти-фрод центр добавляет официальные предупреждения. Не обязательно включать всё на максимум, иначе можно заблокировать курьера, доставку и врача — разрешите звонки из контактов и недавних абонентов.

\end{tcbinfo}

\begin{tcbpractice}{Практика: включите блокировку спама сегодня же}
\begin{enumerate}
  \item Откройте настройки приложения «Телефон», найдите «Блокировка спама / Неизвестные вызовы беззвучно» и включите умную блокировку.
  \item Проверьте телефон родителей: тот же путь + защита оператора (China Mobile — отправьте KTFSR на 10086).
  \item Установите «National Anti-Fraud Center» из магазина приложений, зарегистрируйтесь и включите предупреждения о звонках.
  \item Позвоните на тестовый номер с маркировкой (или посмотрите журнал блокировки), чтобы убедиться, что функция работает.
\end{enumerate}

\end{tcbpractice}

\booksection{Предупреждения о погоде, стихийных бедствиях и землетрясениях}

Блокировка спама защищает от «телефонных мошенников», предупреждения — от «стихий и экстремальной погоды». Многие ставят погодное приложение только ради температуры, но не включают уведомления и предупреждения о землетрясениях. За секунды и десятки секунд до красного штормового предупреждения, тайфуна или землетрясения телефон может дать драгоценное время на спасение.

\booksubsection{Погодные предупреждения: как получать шторм, тайфуны и жару}

\begin{itemize}
  \item Официальные каналы — прежде всего: самые надёжные предупреждения публикуют Китайское метеорологическое управление / Центральная метеорологическая обсерватория / метеослужбы провинций. Можно подписаться на официальный WeChat-аккаунт и Weibo «China Weather» или установить официальное приложение «China Weather».
  \item Системное погодное приложение: Huawei Weather, Xiaomi Weather, OPPO Weather, vivo Weather, «Погода» Apple и др. — большинство присылает предупреждения местной метеослужбы. Откройте настройки погоды и убедитесь, что уведомления «Метеопредупреждения / Стихийная погода» включены.
  \item Минутные прогнозы дождя вроде Caiyun Weather: удобно перед дорогой на работу или за ребёнком — «нужен ли зонт на ближайшие полчаса». Но официальные предупреждения метеослужбы приоритетнее.
  \item Цвета предупреждений (общие): Синий: будьте внимательны, обычно можно ездить как обычно, следите за обновлениями.
  \item Жёлтый: усильте осторожность, сократите ненужные прогулки.
  \item Оранжевый: серьёзно, по возможности оставайтесь в помещении, держитесь подальше от опасных зон.
  \item Красный: особо серьёзно, выполняйте указания о закрытии школ/работ, не проезжайте затопленные участки, ущелья и штормовые дороги.
\end{itemize}

Типы предупреждений: ливни, снегопады, тайфуны, жара, морозы, сильный ветер, грозы, град, гололёд, туман, засуха, риск геологических бедствий.

\booksubsection{Предупреждения о стихийных бедствиях: не только «дождь»}

\begin{itemize}
  \item Кто публикует: метеослужбы — погоду; геологические и селевые риски часто совместно публикуют природоохранные, водные и МЧС-структуры; местные власти используют экстренное радио, SMS, 12379 и другие каналы.
  \item Как получать по многим каналам: Системные уведомления + предупреждения в погодном приложении.
  \item WeChat: подпишитесь на местные аккаунты «Управление по чрезвычайным ситуациям», «Метеослужба».
  \item Каналы 12379 — государственные предупреждения (SMS/сервисы, зависит от региона).
  \item Телевизор, экстренное радио, рупоры в деревнях/сообществах — в горах и в паводковый сезон не полагайтесь только на телефон.
\end{itemize}

На улице и в поездках: при тайфуне не подходите к щитам и деревьям; при ливне не заезжайте в подземные парковки, тоннели и путепроводы; в горах при предупреждении о геологическом риске немедленно уходите со склонов и ущелий.

Связь с памяткой по первой помощи: при потере связи, разряженном телефоне и необходимости вызвать помощь — смотрите главу 32 «Землетрясения и катастрофы», а также главу 7 «Спутниковая связь и связь без сети».

\booksubsection{Предупреждение о землетрясении: это не «прогноз»}

Предупреждение о землетрясении (EEW) — это сигнал, который приходит до сейсмических волн (особенно разрушительных поперечных) в потенциально пострадавший район за секунды и десятки секунд: электромагнитная волна быстрее сейсмической, а станции раньше фиксируют продольную волну. Оно не предсказывает, будет ли землетрясение, но даёт окно «спрятаться рядом, отключить огонь, выйти из лифта». Чем ближе к эпицентру, тем короче предупреждение — иногда его вовремя не хватает.

\begin{itemize}
  \item Официальные и отраслевые решения: в материковом Китае распространены возможности Китайской сети предупреждений о землетрясениях, а также сервисы Чэндуского института сейсмозащиты и других организаций; часть телефонов поставляется с предустановкой или принимает предупреждения онлайн.
  \item Сторонние приложения: ищите в магазине официальные/институциональные приложения «Предупреждение о землетрясении» (проверяйте разработчика, чтобы не нарваться на подделки, требующие кучу лишних прав).
  \item Что делать после сигнала: в помещении — спрячьтесь поблизости (под столом, у несущей стены, рядом с прочной мебелью), защитите голову и шею; держитесь подальше от окон, внешних стен и высокой мебели; не бегите слепо по лестнице и не прыгайте; из лифта выходите на ближайшем этаже; на улице — на открытое место, вдали от проводов, щитов и аварийных зданий. После толчка проверьте газ и проводку, остерегайтесь афтершоков.
  \item Не отключайте и не глушите: предупреждение должно уметь «прорываться» сквозь беззвучный режим и «Не беспокоить»; на ночь кладите телефон под руку и не ставьте громкость на 0.
\end{itemize}

\booksubsection{Встроенное предупреждение о землетрясении: как включить по брендам}

Пути входа могут немного отличаться в разных версиях — ищите в настройках «землетрясение», «предупреждение», «экстренное предупреждение». iOS в материковом Китае поддерживает такие предупреждения ограниченно — обычно нужно установить надёжное стороннее приложение и разрешить геолокацию и уведомления.

\begin{center}
\begin{tabularx}{\linewidth}{|X|X|X|}
\hline
\textbf{Платформа} & \textbf{Частые входы} & \textbf{Пояснение} \tabularnewline
\hline
Xiaomi / Redmi & Менеджер телефона / Центр безопасности > Предупреждение о землетрясении; или поиск в настройках «землетрясение» & В части версий сервис встроен в систему, нужен интернет и включённая геолокация \tabularnewline
Huawei / Honor & Настройки > Безопасность > Экстренные уведомления / Предупреждение о землетрясении & Может называться «Экстренные уведомления», включает землетрясения и др. \tabularnewline
OPPO / realme / OnePlus & Настройки > Безопасность (SOS-связь) > Предупреждение о землетрясении & Зависит от модели и покрытия сервиса в регионе \tabularnewline
vivo / iQOO & Погодное приложение или Настройки > Безопасность/Дополнительно > Предупреждение о землетрясении & Ориентируйтесь на фактическое меню вашей версии системы \tabularnewline
iPhone / iOS & В системных настройках часто нет полного переключателя; установите надёжное приложение предупреждения и разрешите уведомления & Приложению нужны фоновая работа и геолокация; не глушите его сигналы режимом «Не беспокоить» \tabularnewline
Samsung и др. & Настройки > Безопасность и чрезвычайные ситуации / Экстренные оповещения & Охват «экстренных оповещений» зависит от региона и оператора \tabularnewline
\hline
\end{tabularx}
\end{center}

\begin{tcbwarn}{Частые заблуждения о предупреждениях}
\begin{itemize}
  \item Предупреждение $\neq$ прогноз: оно не говорит, «будет ли землетрясение завтра», а срабатывает, когда землетрясение уже началось, но волны ещё не дошли.
  \item Не везде есть запас в десятки секунд: у эпицентра может быть лишь несколько секунд или вовсе не хватить — в основном предупреждение помогает окрестным районам.
  \item Выключили уведомления — всё равно что не ставили: даже лучшее приложение не сработает при беззвучном режиме, убийстве фона и отсутствии геолокации.
  \item Подделки: остерегайтесь «прогнозов землетрясений», запрашивающих доступ к контактам и SMS; предпочитайте встроенные решения или сервисы официальных организаций.
\end{itemize}

\end{tcbwarn}

\begin{tcbinfo}{Как сочетать возможности предупреждений}
Рекомендуется: «Системные погодные уведомления + официальные метео/МЧС-каналы + предупреждение о землетрясении (система или надёжное приложение)». Как и со спамом, не нужно ставить кучу однотипных программ; важнее разрешить уведомления, геолокацию и «прорыв» через «Не беспокоить». После настройки блокировки спама на телефоне родителей проверьте им и погодные предупреждения, и предупреждения о землетрясениях.

\end{tcbinfo}

\begin{tcbpractice}{Практика: проверьте предупреждения за 10 минут}
\begin{enumerate}
  \item Откройте системное погодное приложение, зайдите в настройки и включите уведомления «Метеопредупреждения / Стихийная погода».
  \item В системных настройках найдите «землетрясение» или «предупреждение» и включите предупреждение о землетрясении/экстренное (Android); на iPhone установите надёжное приложение и разрешите уведомления.
  \item Подпишитесь на аккаунты местной метеослужбы или МЧС и разрешите им push.
  \item Убедитесь, что звук предупреждений не полностью глушится режимом «Не беспокоить»; то же самое проверьте на телефоне родителей.
\end{enumerate}

\end{tcbpractice}

\booksection{Межсетевое взаимодействие: удалённый доступ, быстрая передача и безопасность аккаунтов}

Телефон удобен, когда он «работает в команде» с компьютером и планшетом. Здесь — полезные инструменты: удалённый доступ к ПК, ИИ-офис и совместная работа, живые обои, передача по локальной сети, общий буфер обмена, одноразовые коды 2FA, а также связка сканирования бумаг с задачами и расписанием. Системнее о мультиустройствах — в главе 17.

\booksubsection{Управление компьютером с телефона: Sunlogin / ToDesk}

\begin{itemize}
  \item Sunlogin Remote Control. Для кого: для тех, кому нужно поработать с домашним или рабочим Windows/Mac, находясь вне дома; также часто используется, чтобы удалённо помочь старшим родственникам с телефоном (см. главу 27). Главные плюсы: зрелая китайская экосистема; с телефона можно управлять рабочим столом компьютера, передавать файлы, включать и выключать его (частично требуется дополнительное оборудование); бесплатной версии обычно хватает для личного нетребовательного использования.
  \item ToDesk Для кого: тем, кому нужны стабильное соединение и низкая задержка для удалённой работы. Ключевые плюсы: подключение по коду устройства + паролю, телефон управляет ПК, ПК — телефоном (зависит от версии), удобная передача файлов и переключение мониторов; в Китае узлы обычно работают хорошо.
  \item Базовое использование: установите клиент на ПК и войдите/держите онлайн установите приложение на телефон создайте сеанс по идентификатору устройства + паролю подключения (или одному аккаунту) управляйте рабочим столом касаниями и жестами.
  \item Красные линии безопасности (запомните обязательно): Если незнакомец, назвавшись «службой поддержки», «полицией, прокуратурой или судом» либо «оператором связи», просит установить Sunlogin или ToDesk и получить удалённый доступ к вашему телефону или компьютеру — отказывайте сразу: это типичная схема мошенничества (подробнее в главе 27).
  \item Пользуясь сами: включите двухфакторную защиту, смените пароль по умолчанию, выходите из аккаунта или отключайте удалённый доступ без присмотра; на рабочем ПК соблюдайте корпоративные правила.
  \item В общественных местах следите, чтобы никто не подглядывал экран и пароли.
\end{itemize}

\booksubsection{ИИ-офис: Trae / WorkBuddy / Dumate}

\begin{itemize}
  \item Trae Для кого: программистам и тем, кто работает с документами; дружелюбен к китайской среде. Ключевые плюсы: ИИ IDE/ассистент от ByteDance, режим SOLO генерирует проект по одной фразе; ускоряет написание кода и правку документов. Если он работает на ПК, с телефона можно следить за задачами через удалённый доступ или официальные мобильные функции (актуальная форма — у продукта). Подробнее — в главе 41 «Панорама ИИ-инструментов для программирования».
  \item WorkBuddy Для кого: офисным работникам, которым ИИ помогает с протоколами встреч, письмами, отчётами и документами. Ключевые плюсы: ИИ-помощник для офисных сценариев, хорош для частых задач «черновик — шлифовка — резюме»; часть процессов можно запустить с телефона, тяжёлую работу оставить на ПК.
  \item Dumate Для кого: командам и отдельным пользователям, которым нужна совместная работа, документы и управление знаниями в одном месте. Ключевые плюсы: ИИ-инструмент для совместной работы с акцентом на сбор материалов и совместное редактирование; на телефоне — просмотр и лёгкие правки, сложную вёрстку лучше на большом экране ПК.
  \item Как понимать «управление компьютером с телефона»: многие ИИ-функции для офиса работают на компьютере или в веб-версии; подключившись с телефона к компьютеру через Sunlogin/ToDesk или воспользовавшись собственным мобильным приложением/веб-версией продукта, вы сможете «дистанционно руководить» ИИ-задачами на компьютере по дороге на работу. Принцип: тяжёлые вычисления и длинные документы — на компьютере; посмотреть прогресс или поправить пару фраз — с телефона.
  \item О приватности: прежде чем загружать договоры, паспорта и данные клиентов, читайте политику продукта и требования компании; для корпоративных материалов используйте инструменты, указанные работодателем.
\end{itemize}

\booksubsection{Живые обои: динамические обои miHoYo (телефон + ПК)}

\begin{itemize}
  \item YoMate. Для кого: для поклонников персонажей miHoYo, желающих бесплатные живые обои. Главные плюсы: бесплатное приложение с живыми обоями от miHoYo, есть версии для телефона и ПК; Live2D-персонажи и анимированные сцены могут стать рабочим столом, а на части моделей и систем — экраном блокировки или обоями.
  \item Экономия и производительность: живые обои едят больше батареи и GPU, чем статичные; при низком заряде или на старых устройствах включайте только при зарядке или снижайте частоту кадров (если приложение позволяет).
  \item Источник установки: предпочитайте официальный сайт/официальный магазин, избегайте «взломанных» версий с троянами; на ПК проверяйте путь установки и необходимость автозапуска.
\end{itemize}

\booksubsection{Быстрая передача файлов по локальной сети: LocalSend / Feige / AirDrop}

\begin{itemize}
  \item LocalSend Для кого: тем, кто смешивает Android/Windows/iOS/Mac/Linux и хочет кроссплатформенную передачу без облака. Ключевые плюсы: open source, прямое соединение в локальной сети, обмен фото/видео/документами в одном Wi-Fi без кабеля и серверов производителя; скорость зависит от роутера.
  \item Feige. Для кого: для офисов с компьютерами Windows и старых офисных локальных сетей. Главные плюсы: проверенная временем утилита передачи и чата по локальной сети, удобна для пересылки больших файлов внутри интранета. Внимание: скачивайте только из официальных источников, а в корпоративной сети соблюдайте правила IT-отдела.
  \item AirDrop Для кого: пользователям iPhone, iPad и Mac. Ключевые плюсы: обмен в один клик в экосистеме Apple, поиск по Bluetooth + прямая передача по Wi-Fi, без трафика; можно ограничить «только контакты», чтобы снизить беспокойство от незнакомцев.
  \item Дополнительно: Xiaomi Mi Share, Huawei Share, vivo EasyShare и другие бренды очень быстры между устройствами одного бренда; «Файловый помощник» WeChat удобен, но ограничен размером и устройствами входа — для больших файлов лучше LocalSend/AirDrop.
  \item Советы: подключите оба устройства к одному Wi-Fi; на публичном Wi-Fi передавать конфиденциальные файлы рискованно — предпочитайте личную точку доступа; после передачи проверьте целостность файла на приёмнике.
\end{itemize}

\booksubsection{Синхронизация буфера обмена: KDE Connect / обмен телефон-ПК}

\begin{itemize}
  \item «Универсальный буфер» Apple Для кого: пользователям iPhone + Mac. Ключевые плюсы: войдите в один Apple ID и включите Handoff — скопировали на телефоне, вставили на Mac (и наоборот); работает с текстом и частью файлов.
  \item KDE Connect Для кого: Android + Windows/Linux/Mac, нужен open source буфер обмена, зеркалирование уведомлений и удалённый ввод. Ключевые плюсы: после сопряжения в локальной сети — синхронизация буфера, пересылка уведомлений с телефона на ПК, телефон как пульт для слайдов/тачпад; данные не проходят через сторонние аккаунты.
  \item Брендовые «обмены телефон-ПК»: Huawei Multi-screen Collaboration / Super Device, Xiaomi MiMo+, vivo EasyShare, OPPO Connect, Honor Connect и др. — часто с межэкранным буфером и перетаскиванием файлов; лучше всего работают внутри одного бренда.
  \item Официальный канал Windows + Android: «Phone Link / Link to Windows» (развитие Your Phone) — на ПК видны уведомления и фото телефона, можно копировать-вставлять в поддерживаемых сценариях; на Samsung и др. интеграция глубже.
  \item О безопасности: в буфере могут остаться пароли и коды; на общем компьютере аккуратнее с глобальной синхронизацией, после использования не вставляйте повторно в поля паролей; часть клавиатур тоже читает буфер — следите за правами.
\end{itemize}

\booksubsection{Одноразовые коды 2FA: Aegis / Microsoft Authenticator}

\begin{itemize}
  \item Почему аутентификатор, а не SMS: SMS-коды могут быть украдены через перевыпуск SIM; TOTP-коды генерируются локально по времени и не зависят от сигнала, устойчивее к фишингу и перехвату (но всё равно нужны защита от троянов и скриншотов).
  \item Aegis Authenticator Для кого: пользователям Android, предпочитающим open source и локальный контроль. Ключевые плюсы: бесплатно и с открытым кодом, шифрованные резервные копии, импорт/экспорт, чистый интерфейс, разблокировка списка токенов отпечатком/лицом.
  \item Microsoft Authenticator Для кого: пользователям аккаунтов Microsoft, Office 365 и корпоративных учёток; работает на iOS и Android. Ключевые плюсы: push-подтверждение входа, мониторинг паролей и резервное копирование аккаунтов (зависит от версии), тесная связь с Windows и корпоративной идентификацией.
  \item Google Authenticator / 2FAS и др.: тоже распространены; главное — наличие шифрованного резервного копирования, чтобы при смене телефона не потерять токены.
  \item При смене телефона обязательно: перенесите со старого или используйте «резервные коды/ключ восстановления» из сервиса; распечатайте резервные коды на бумаге и храните в двух местах, не только на одном телефоне.
  \item С менеджерами паролей: часть хранилищ умеют хранить TOTP — удобно, но риск сосредоточен в одном месте; для ценных аккаунтов лучше разделить «менеджер паролей + отдельный аутентификатор». Подробнее — в главе 14.
\end{itemize}

\booksubsection{Рекомендуемые связки: как выстроить день}

\textbf{Пример рабочего процесса}\par

\begin{itemize}
  \item Утро: посмотреть расписание в TickTick глянуть погоду/предупреждения открыть Trae/документы на ПК, телефон отложить.
  \item День: бумаги сканировать в PDF и отправить на ПК через LocalSend/AirDrop; длинные тексты чертить на ПК с ИИ-офисными инструментами.
  \item Вне дома: подключайтесь к компьютеру через ToDesk/Sunlogin и продолжайте свои задачи; вход в важные аккаунты — через одноразовые коды Aegis/Microsoft Authenticator.
  \item Вечер: отмечать выполненное в TickTick; при желании сменить обои в «Живых обоях» и расслабиться.
\end{itemize}

\begin{tcbwarn}{Правила удалённого доступа и безопасности аккаунтов}
\begin{itemize}
  \item Кто-то просит установить удалённый доступ положите трубку и сами найдите официальную поддержку.
  \item Коды подтверждения и одноразовые пароли не сообщайте никому, включая «курьера» и «службу поддержки».
  \item Перед удалённым доступом к рабочему ПК и установкой ИИ-программ спросите разрешения у IT-отдела.
  \item В публичных сетях будьте особенно внимательны при передаче файлов и удалённом входе.
\end{itemize}

\end{tcbwarn}

\begin{tcbpractice}{Практика: прогоните цепочку «между устройствами» один раз}
\begin{enumerate}
  \item Подключите телефон и ПК к одному Wi-Fi и передайте фото через LocalSend (или AirDrop).
  \item Создайте в TickTick задачу «Сегодня в 20:00 разобрать сканы» и убедитесь, что она видна на телефоне и ПК.
  \item Включите 2FA для почты или WeChat (если поддерживается), привяжите Aegis или Microsoft Authenticator и сохраните резервные коды.
  \item Установите Sunlogin или ToDesk на собственный компьютер, один раз успешно подключитесь со своего телефона, затем выйдите из программы или усильте пароль — вы освоите процесс, не оставив «чёрного хода» в безопасности.
\end{enumerate}

\end{tcbpractice}
